% Abstract (<= 200 words). Sources: core/universe_counts.csv; features/fig_name_share_by_year.csv;
% d1_break/RESULTS.md (year-dummy deviations; H1b 0.37-0.65; Shapley 0.574; H1e 7.1%, 16.3%, 33.3%);
% d2_carriers/RESULTS.md (H2d); d3_register/RESULTS.md (H3a); d4_filter/RESULTS.md (H4a, H4e ratio 0.50).
\begin{abstract}
At ICLR, machine-learning paper titles have turned into brands. We study 43,482 ICLR submissions from 2017 to 2026, the pre-specified discovery split (about 80\%) of the submission record, including rejected, withdrawn and desk-rejected papers. Titles of the form ``NAME: subtitle'' rose from 8.9\% to 40.5\% of submissions on a steady logistic path: 2024 and 2025 are consistent with the 2017--23 trend, and 2026 lies 5.5 points above it. Depending on design, 37--65\% of the 2023--26 rise tracks name-masked abstract content. Team composition adds little beyond content. Detectable LLM drafting accounts for 7\% under the pre-specified fingerprint and 16--33\% under broader markers. In the one additive split, 57\% remains unexplained. The rise is graded by how new a team's profile keys are, and incumbent keys, pooled over list positions, take up new title style more slowly than new topics. LLM-favoured verbs were added to meaning-fixed title openers from 2024. Accepted papers carry named titles at the all-submission rate, also after correcting for camera-ready edits. Published papers alone estimate about half the acceleration, imprecisely (ratio 0.50, 95\% CI $-0.15$ to 1.16). We report failed predictions from a locked, pilot-informed plan, a reusable camera-ready audit, and cautions for studying co-authored text.
\end{abstract}
